% !PRINTED 556
\weberpage{556}
%<<CONT
absolut genommen nicht grösser als $k$. Es giebt eine durch $\mu$ theilbare von Null verschiedene Zahl:

\weq{5}{\alpha = a_1 \omega_1 + a_2 \omega_2 + \cdots + a_n \omega_n,}

in der die ganzzahligen Coëfficienten $a_1, a_2 \ldots a_n$ die Grenzen $\pm k$ nicht überschreiten, und folglich ist nach (2) und (3)

\weq{6}{N_a(\alpha) < R k^n < R N_a(\mu).}

Da nun $\alpha$ durch $\mu$ theilbar ist, so setzen wir

\weq{7}{\alpha = \mu\varphi,\ N_a(\alpha) = N_a(\mu) N_a(\varphi),}

und erhalten aus (6)

\weq{8}{N_a(\varphi) < R.}

Wenn nun $\psi$ ein Repräsentant einer beliebig gegebenen Classe $C$ ist, so wählen wir $\mu$ so, dass

\[
\beta = \mu\psi
\]

eine Zahl ist, und wenn dann nach (7)

\[
\alpha = \mu\varphi
\]

ist, so ist

\[
\frac{\varphi}{\psi} = \frac{\alpha}{\beta},
\]

also $\varphi$ und $\psi$ äquivalent. $\varphi$ ist also gleichfalls ein Repräsentant der Classe $C$, und dieser genügt der Bedingung (8). Es kommt also, wie bewiesen werden sollte, in jeder Classe ein Functional vor, dessen absolute Norm unter $R$ liegt.

Die Anzahl der Idealclassen, die wir mit $h$ bezeichnen wollen, ist hiernach eine dem Körper $\Omega$ eigenthümliche natürliche Zahl, die die \gesperrt{Classenzahl} genannt wird.

In dem einfachsten Falle, wo die Classenzahl gleich 1 ist, ist jedes Functional mit einer Zahl associirt, d. h. es kann jedes (ganze oder gebrochene) Functional durch Absonderung eines Zahlenfactors in eine Einheit verwandelt werden. In diesem Falle lässt sich jede ganze Zahl des Körpers $\Omega$ in \gesperrt{Primzahlfactoren} zerlegen, und diese Körper haben eine Theorie, die im Wesentlichen mit der rationalen Zahlentheorie übereinstimmt. Für solche Körper ist die Einführung der Functionale und Ideale nicht nothwendig.

Hierher gehören neben dem Körper der rationalen Zahlen unter anderen der Körper der \gesperrt{Gauss}'schen imaginären Zahlen und der aus dritten Einheitswurzeln gebildeten Zahlen (Bd. I, \S.\,173, 174).
